\begin{lemma}[All ordered one-point extensions remain in the tree]
Let $F$ be a front on the full base $\mathbb N$.  If
$s\in T(F)\setminus F$ and $n$ is such that
$\mathsf{ProperInitialSegment}(s,s\cup\{n\})$, then
$s\cup\{n\}\in T(F)$.
\end{lemma}
\begin{proof}
Choose $t\in F$ with $\mathsf{InitialSegment}(s,t)$ using the definition of
\noderef{PrefixTree}.  Since $t\in F$ and $s\notin F$, we have $s\ne t$.
By \noderef{InitialSegment}, choose $r\in t$ such that
$s=\{k\in t\mid k<r\}$.  Thus $s$ is a proper initial segment of $t$ in
the sense of \noderef{ProperInitialSegment}.

Put $v=s\cup\{n\}$.  We now extract the precise order information from
the hypothesis $\mathsf{ProperInitialSegment}(s,v)$.  By
\noderef{ProperInitialSegment} and \noderef{InitialSegment}, the equality
case $s=v$ is impossible, so there is a cutoff $c\in v$ such that
\[
  s=\{k\in v\mid k<c\}.
\]
The cutoff $c$ cannot belong to $s$: if $c\in s$, the displayed equality
would imply $c<c$.  Hence $c\in v\setminus s$.  Since
$v=s\cup\{n\}$, this gives $c=n$ and $n\notin s$.  The displayed equality
therefore says, explicitly, that every $k\in s$ satisfies $k<n$ and that
every element of $v$ below $n$ belongs to $s$.
In particular, $\max(v)=n$.

Consider the infinite set
\[
  Y=v\cup\{k\in\mathbb N\mid \max(v)<k\}.
\]
Because $\max(v)=n$, the set $Y\setminus v$ consists exactly of the
naturals greater than $n$.  Thus $Y$ contains all naturals greater than
$n$, so it is infinite, and clearly $Y\subseteq\mathbb N$.  Moreover, the
cutoff equality above becomes
\[
  s=\{k\in Y\mid k<n\},
\]
because the elements of $Y\setminus v$ are all greater than $n$.  Density
in \noderef{Front} now gives $w\in F$ and
$\mathsf{ProperPrefixSet}(w,Y)$, using the definition in
\noderef{ProperPrefixSet}.  Let $p\in Y$ be its cutoff, so
\[
  w=\{k\in Y\mid k<p\}.
\]
Compare the natural numbers $p$ and $n$.
If $p=n$, then $w=s$, contradicting $w\in F$ and $s\notin F$.  If
$p<n$, then $p\in s$ and the two cutoff descriptions show directly that
$w$ is a proper initial segment of $s$.  Moreover, $p\in s$ implies
$p<r$, and every $k\in t$ with $k<p$ belongs to $s$ because $p<r$.
Consequently
$w=\{k\in t\mid k<p\}$, so $\mathsf{InitialSegment}(w,t)$ by
\noderef{InitialSegment}.  Since both $w$ and $t$ belong to $F$,
prefix-freeness in \noderef{Front} gives $w=t$.  But $w$ is a proper
initial segment of $s$, while $s$ is a proper initial segment of $t$, so
$w\ne t$, a contradiction.  Consequently $n<p$.

It follows from the formula for $w$ that $s\subseteq w$ and $n\in w$,
so $v\subseteq w$.  There are now two cases.  If $w=v$, then
$w\in F$ and $\mathsf{InitialSegment}(v,w)$ holds by equality; hence
$v\in\mathsf{PrefixTree}(F)$ by \noderef{PrefixTree}.  Suppose instead
that $w\ne v$.  Choose the least element $q$ of the nonempty finite set
$w\setminus v$.  Every element of $v$ is at most $n$, while every element
of $Y\setminus v$ is greater than $n$; hence $n<q$.  If $k\in w$ and
$k<q$, then $k\in v$, by the minimality of $q$.  Conversely every
$k\in v$ satisfies $k<q$, because the elements of $s$ are below $n$ and
$n<q$.  Therefore
\[
  v=\{k\in w\mid k<q\},
\]
with $q\in w$, so \noderef{InitialSegment} gives
$\mathsf{InitialSegment}(v,w)$.  Since $w\in F$, the definition of
\noderef{PrefixTree} now yields $v\in\mathsf{PrefixTree}(F)$.
This argument is exactly the ordered-extension form of the front-density
step used at paper lines 1145--1148.
\end{proof}
